\documentclass[10pt]{article}
\usepackage[a4paper,margin=2.1cm]{geometry}
\usepackage{amsmath,amssymb,booktabs,enumitem}
\usepackage[T1]{fontenc}\usepackage{lmodern}
\setlength{\parindent}{0pt}\setlength{\parskip}{5pt}
\newcommand{\bs}{\beta^\star}\newcommand{\E}{\mathbb E}\newcommand{\asinh}{\operatorname{arcsinh}}
\newcommand{\step}[1]{\par\smallskip\textbf{Step #1.}\ }
\title{Derivations: squared-loss unlearning in the diagonal network}
\author{}\date{25 November 2026}
\begin{document}\maketitle

Sections 1--7 are exact (finite $d$, any data). Sections 8--11 derive the large-$d$ end points with the replica method, step by step; Section 12 rederives them with the cavity method as a check; Sections 13--14 give the trajectories and the solver. Replica and cavity are derivations, not proofs. Each numbered result is followed by its derivation.

\section{Setting}
Network: $f(x)=x^\top\beta$, $\beta=u^2-v^2$ entrywise, $u,v\in\mathbb R^d$, positive. The samples $x_i$ are the columns of $X$. A training stage minimises
\[
L=\tfrac12\sum_i w_i\,(x_i^\top\beta-t_i)^2
\]
by gradient flow on $(u,v)$, with weights $w_i$ and targets $t_i$ (pretraining: $w_i=1$, $t_i=y_i$). Write $G=\nabla_\beta L=\sum_iw_i(x_i^\top\beta-t_i)\,x_i\in\mathbb R^d$.

\section{The flow in one variable}
\textbf{Result 1.} Along gradient flow, $c=2uv$ is constant in time (entrywise), and
\begin{equation}\label{mf}
\beta=c\sinh\theta,\qquad\dot\theta=-4G,\qquad\theta=\log(u/v).
\end{equation}
\step1 Chain rule. $\partial\beta_j/\partial u_j=2u_j$ and $\partial\beta_j/\partial v_j=-2v_j$, so
\[
\dot u=-\frac{\partial L}{\partial u}=-2u\,G,\qquad\dot v=-\frac{\partial L}{\partial v}=+2v\,G\quad\text{(entrywise).}
\]
\step2 The product is conserved: $\frac{d}{dt}(uv)=\dot uv+u\dot v=-2uvG+2uvG=0$.
\step3 The ratio moves linearly in $G$: $\frac{d}{dt}\log u=-2G$, $\frac{d}{dt}\log v=+2G$, so $\dot\theta=-4G$.
\step4 Express $\beta$. From $uv=c/2$ and $u/v=e^\theta$: $u^2=\frac c2e^{\theta}$, $v^2=\frac c2e^{-\theta}$, so $\beta=\frac c2(e^\theta-e^{-\theta})=c\sinh\theta$.

\textbf{Consequence.} $G$ is a combination of the $x_i$, so $\theta(t)-\theta(0)\in\operatorname{span}X$ at all times. Also $u^2+v^2=c\cosh\theta=\sqrt{\beta^2+c^2}$, hence
\begin{equation}\label{metric}
\dot\beta=c\cosh\theta\,\dot\theta=-4\,(u^2+v^2)\,G .
\end{equation}
The squared weight norm $u_j^2+v_j^2$ of coordinate $j$ is the rate at which it moves.

\section{Where a stage ends}
Define the convex function $\psi_c$ by $\psi_c'(\beta)=\asinh(\beta/c)$, $\psi_c(0)=0$:
\[
\psi_c(\beta)=\int_0^\beta\asinh(s/c)\,ds=\beta\,\asinh(\beta/c)-\sqrt{\beta^2+c^2}+c ,
\]
and its Bregman divergence $D_c(\beta,\beta_0)=\psi_c(\beta)-\psi_c(\beta_0)-\psi_c'(\beta_0)(\beta-\beta_0)$, summed over coordinates.

\textbf{Result 2.} If a stage starts at $\beta_0$ with scale $c$ and converges to zero loss on its samples, it ends at
\begin{equation}\label{proj}
\beta_{\rm end}=\arg\min_\beta D_c(\beta,\beta_0)\quad\text{s.t.}\quad x_i^\top\beta=t_i\ \ \forall i .
\end{equation}
\step1 Write the optimality conditions of \eqref{proj}. With one multiplier $a_i$ per constraint, stationarity in $\beta$ reads $\psi_c'(\beta)-\psi_c'(\beta_0)=-Xa$, i.e.\ $\asinh(\beta/c)-\asinh(\beta_0/c)\in\operatorname{span}X$, together with the constraints.
\step2 The flow satisfies them. By Result 1, $\asinh(\beta/c)=\theta$ and $\theta_{\rm end}-\theta_0=-4\int_0^\infty G\,dt=-Xa$ with $a_i=4\int_0^\infty w_i(x_i^\top\beta-t_i)dt$. Zero loss gives the constraints.
\step3 $\psi_c$ is strictly convex, so the point satisfying these conditions is unique.

\textbf{Result 3 (two regimes).} Per coordinate:
\begin{align*}
|\beta|,|\beta_0|\gg c,\ \text{same sign:}&\quad D_c\approx|\beta|\log\frac{|\beta|}{|\beta_0|}-|\beta|+|\beta_0| &&\text{(does not depend on $c$)},\\
\beta_0=0,\ |\beta|\gg c:&\quad D_c\approx|\beta|\Big(\log\frac{2|\beta|}{c}-1\Big) &&\text{(a new coordinate costs }\propto\log\tfrac1c),\\
|\beta|,|\beta_0|\ll c:&\quad D_c\approx\frac{(\beta-\beta_0)^2}{2c} &&\text{(squared distance)}.
\end{align*}
\step1 For $|s|\gg c$, $\asinh(s/c)\approx\operatorname{sign}(s)\log(2|s|/c)$ and $\sqrt{s^2+c^2}\approx|s|$, so $\psi_c(s)\approx|s|(\log\frac{2|s|}c-1)$. This is the second line.
\step2 Insert into $D_c$ for $\beta,\beta_0>0$: $\beta(\log\frac{2\beta}c-1)-\beta_0(\log\frac{2\beta_0}c-1)-\log\frac{2\beta_0}c(\beta-\beta_0)=\beta\log\frac\beta{\beta_0}-\beta+\beta_0$.
\step3 For $|s|\ll c$, $\asinh(s/c)\approx s/c$, so $\psi_c(s)\approx s^2/2c$ and $D_c$ is its quadratic Bregman divergence.

The first two lines are the \emph{rich} regime (changing existing weights is cheap, creating new ones is expensive); the third is the \emph{lazy} regime.

\section{What each operation does}
\textbf{Result 4 (rescaling).} Multiplying $u,v$ by $\sqrt S$ gives $\beta\to S\beta$ and $c\to Sc$. Finetuning after it ends at
\begin{equation}\label{resc}
\beta_{\rm end}=S\cdot\arg\min_{\beta'}D_c(\beta',\beta_{\rm pre})\quad\text{s.t.}\quad x_i^\top\beta'=t_i/S .
\end{equation}
In words: rescaling by $S$ is the same as finetuning the unrescaled model to targets divided by $S$, then multiplying the result by $S$.
\step1 $\beta=u^2-v^2\to S\beta$ and $c=2uv\to Sc$.
\step2 By Result 2 the end point minimises $D_{Sc}(\beta,S\beta_{\rm pre})$ subject to $x_i^\top\beta=t_i$.
\step3 $\psi_c(\beta)=c\,\phi(\beta/c)$ with $\phi(x)=x\asinh x-\sqrt{x^2+1}+1$, so $\psi_{Sc}(S\beta')=S\psi_c(\beta')$ and $\psi_{Sc}'(S\beta')=\psi_c'(\beta')$. Hence $D_{Sc}(S\beta',S\beta_{\rm pre})=S\,D_c(\beta',\beta_{\rm pre})$. Substitute $\beta=S\beta'$.

\textbf{Corollary (no shared block, forget-to-zero).} The forget block ends at $S$ times where it ends for $S=1$.
\step1 With no shared block the retain samples have no entries on the forget block, so the constraints, and therefore \eqref{resc}, separate into a retain problem and a forget problem.
\step2 The forget problem has targets $0/S=0$: its solution $\beta'_f$ does not depend on $S$, and $\beta_{f,\rm end}=S\beta'_f$.

Consequences, with $\delta_k=x_k^\top\beta'_f-y_k$ the change forget-to-zero makes on a left-over sample at $S=1$: its error after rescaling is $((S-1)y_k+S\delta_k)^2$, and the distance to the retain-only model (whose forget block is zero, Result 8) is $S^2\|\beta'_f\|^2/d$. $S\to0$ erases the class; $S>1$ re-amplifies what forgetting removed, on every direction except $F$ itself, which stays pinned at zero. Simulation, blocks $(1,0,1)$: distance $0.002,0.002,0.012,0.096$ at $S=0.1,0.3,1,3$; $S^2\times0.012$ gives $0.0001,0.001,0.012,0.108$.

\textbf{Result 5 (weight decay).} With the penalty $\frac\lambda2(\|u\|^2+\|v\|^2)$ added to $L$: $c(t)=c(0)e^{-2\lambda t}$ and $\dot\theta=-4G$ as before. If it stays on until convergence, the end point solves
\[
\min_\beta\ L(\beta)+\tfrac\lambda2\|\beta\|_1 .
\]
\step1 The penalty adds $-\lambda u$ and $-\lambda v$: $\dot u=-2uG-\lambda u$, $\dot v=2vG-\lambda v$.
\step2 $\frac d{dt}(uv)=-2\lambda uv$, so $c$ decays; $\frac d{dt}\log(u/v)=-4G$, unchanged.
\step3 Stationary points: for each coordinate, $u(2G+\lambda)=0$ and $v(2G-\lambda)=0$. Since $c\to0$, at most one of $u,v$ survives. If $\beta>0$ ($u>0$): $G=-\lambda/2$. If $\beta<0$: $G=+\lambda/2$. If $\beta=0$: $|G|\le\lambda/2$ (otherwise one of $u,v$ grows). These are the optimality conditions of the lasso problem above.

\textbf{Result 6 (SGD).} One step of size $\gamma$ with mini-batch gradient $G$ is exactly
\[
\theta\leftarrow\theta-2\operatorname{artanh}(2\gamma G),\qquad c\leftarrow c\,(1-4\gamma^2G^2).
\]
\step1 The step is $u\leftarrow u(1-2\gamma G)$, $v\leftarrow v(1+2\gamma G)$.
\step2 $uv\leftarrow uv(1-4\gamma^2G^2)$; $\log\frac uv\leftarrow\log\frac uv+\log\frac{1-2\gamma G}{1+2\gamma G}=\log\frac uv-2\operatorname{artanh}(2\gamma G)$.

So weight decay and SGD can only shrink $c$.

\section{How forgetting spreads: the local rule}
\textbf{Result 7.} To first order, finetuning from $\beta_0$ to targets $t$ on samples $X$ changes $\beta$ by
\begin{equation}\label{lin}
\Delta\beta=M\,X\,(X^\top MX)^{-1}(t-X^\top\beta_0),\qquad M=\operatorname{diag}(u_j^2+v_j^2)\ \text{at }\beta_0 .
\end{equation}
\step1 By \eqref{metric}, $\dot\beta=-4MG$ with $G\in\operatorname{span}X$. Freezing $M$ at its initial value, $\Delta\beta=MX\lambda$ for some $\lambda$.
\step2 Impose the constraints $X^\top(\beta_0+\Delta\beta)=t$ and solve for $\lambda$.

\textbf{Spread to an input not in the loss.} Take forget-to-zero on $n_F$ samples and ignore the retain term. Let $x$ be a forget-class input independent of $\Delta\beta$, and $m_F$ the mean squared change on $F$. Then
\begin{equation}\label{spread}
\frac{\E(x^\top\Delta\beta)^2}{m_F}\approx\frac{n_F}{d_{\rm eff}},\qquad d_{\rm eff}=\frac{(\sum_jM_j)^2}{\sum_jM_j^2}\quad\text{(sums over the coordinates the forget class sees).}
\end{equation}
\step1 $x$ has independent entries of variance $1/d$ on those coordinates, so $\E(x^\top\Delta\beta)^2=\|\Delta\beta\|^2/d$.
\step2 With $r=t-X^\top\beta_0$: $\|\Delta\beta\|^2=r^\top K^{-1}(X^\top M^2X)K^{-1}r$, $K=X^\top MX$.
\step3 For $n_F\ll d_{\rm eff}$ the samples are nearly orthogonal in these metrics: $K\approx\frac1d\sum_jM_j\,I$ and $X^\top M^2X\approx\frac1d\sum_jM_j^2\,I$. Insert, and use $\|r\|^2=n_Fm_F$.

Flat metric ($M=I$): $d_{\rm eff}$ is the number of coordinates the class sees. After small-$k$ pretraining $M_j=\sqrt{\beta_j^2+c^2}$ is concentrated on the learned coordinates and $d_{\rm eff}$ is their effective number. Baseline: $d_{\rm eff}=0.28$ of the class coordinates. Error on left-over samples relative to error on $F$, at $|F|/|F_{\rm full}|=0.1,0.5,0.9$: \eqref{spread} gives $0.11,0.53,0.96$; \eqref{lin} with the retain term $0.13,0.55,0.77$; simulation $0.12,0.40,0.72$; flat metric $0.03,0.16,0.33$. (For a left-over sample $x$ is not independent of $M$; this is neglected.) Combined with Result 4 for the rescaling: $0.85,0.77,0.55,0.43$ at $S=0.1,0.3,1,3$ against $0.67,0.56,0.40,0.35$ in simulation.

\section{Facts about losses and references}
\textbf{Result 8.} The retain-only model is zero on the forget block. \emph{Derivation:} it starts at $\theta=0$ and $\theta$ stays in the span of the retain samples (Result 1), which are zero on the forget block; $\beta=c\sinh0=0$ there.

\textbf{Result 9.} With no shared block a forget term cannot change the retain block. \emph{Derivation:} its gradient is a combination of forget samples, zero on the retain block; by \eqref{mf} $\theta$ does not move there.

\textbf{Result 10.} From a model that fits $F$, the signed loss and pure ascent do not move. \emph{Derivation:} their forget gradient is $-\mu_f\sum_F(x_i^\top\beta-y_i)x_i=0$, and the retain gradient is zero as well.

\textbf{Result 11.} A target $b\,y_i$ fitted on $F$ leaves error $(1-b)^2\,\overline{y^2}$ on $F$. \emph{Derivation:} $x_i^\top\beta-y_i=(b-1)y_i$. (Simulation: $1.311$, $2.341$ for $b=-0.5,-1$; formula $1.31$, $2.33$.)

\section{Untied layers (Azulay et al.\ 2021)}
Network per coordinate: $w=u_+v_+-u_-v_-$. Let $g=\nabla_wL$ and $R=\frac12(u_+^2+v_+^2+u_-^2+v_-^2)$.

\textbf{Result 12.} $\dot w=-2R\,g$; $c^2=R^2-w^2$ is conserved; $w=c\sinh\theta$ with $\dot\theta=-2g$. Results 2--4 and 7 hold with this $c$ and $M=2R$.
\step1 $\dot u_+=-gv_+$, $\dot v_+=-gu_+$, $\dot u_-=+gv_-$, $\dot v_-=+gu_-$.
\step2 $\dot w=\dot u_+v_++u_+\dot v_+-\dot u_-v_--u_-\dot v_-=-g(u_+^2+v_+^2+u_-^2+v_-^2)=-2Rg$.
\step3 $\dot R=u_+\dot u_++v_+\dot v_++u_-\dot u_-+v_-\dot v_-=-2g\,u_+v_++2g\,u_-v_-=-2g\,w$.
\step4 $\frac d{dt}(R^2-w^2)=2R(-2gw)-2w(-2Rg)=0$.
\step5 With $c$ constant and $R=\sqrt{w^2+c^2}$: $\dot w=-2g\sqrt{w^2+c^2}$, i.e.\ $\frac d{dt}\asinh(w/c)=-2g$.

Special cases. $u_+=u_-=u_0$, $v_+=v_-=v_0$ at the start: $w=0$, $c=u_0^2+v_0^2$, which is their Theorem 4.1 ($\sqrt k=2c$). $u_\pm=v_\pm$: our network, $c=2u_+u_-$. Steps 1--5 are my derivation for a general start, checked numerically; their general statement (Appendix B) was not in the copy I read.

\textbf{Layer rescaling.} $u_\pm\to au_\pm$, $v_\pm\to v_\pm/a$ keeps $w$. On a trained coordinate of our network ($u_+=v_+=\sqrt w$, $u_-=v_-\approx0$): $R=\frac w2(a^2+a^{-2})$, so $c=\sqrt{R^2-w^2}=\frac w2|a^2-a^{-2}|$. It does not exist for $\beta=u^2-v^2$.


\section{Large $d$, replica method: one stage}\label{rep1}
From here on $x_i\sim\mathcal N(0,I/d)$, $n=\alpha d$ samples, $d\to\infty$. We compute where one stage ends when its start does not depend on the data.

\textbf{Problem.} By Result 2,
\[
\hat\beta=\arg\min_\beta\sum_j\phi_j(\beta_j)\quad\text{s.t.}\quad x_i^\top\beta=t_i,\qquad\phi_j(\beta)=D_c(\beta,\beta_{0,j}),\qquad t_i=x_i^\top b+\sigma\xi_i ,
\]
with $b\in\mathbb R^d$ the vector that generates the targets (pretraining: $b=\bs$, $\beta_0=0$), $\xi_i\sim\mathcal N(0,1)$ label noise, and $(\beta_{0,j},b_j)$ independent across $j$ and of $X$. Write $w=\beta-b$, so the constraints are $x_i^\top w=\sigma\xi_i$.

\textbf{Result 13.} Each coordinate of $\hat\beta$ is distributed as the solution of
\begin{equation}\label{ch1}
\phi'(\hat\beta)+A(\hat\beta-b)=\eta,\qquad\eta\sim\mathcal N(0,r),
\end{equation}
\begin{equation}\label{fp1}
A=\frac\alpha\chi,\qquad r=\frac{\alpha\,(q+\sigma^2)}{\chi^2},\qquad q=\E(\hat\beta-b)^2,\qquad\chi=\E\frac{\partial\hat\beta}{\partial\eta}=\E\frac1{\phi''(\hat\beta)+A}.
\end{equation}
Two unknown numbers, $q$ and $\chi$. Here $\phi'(\beta)=\asinh(\beta/c)-\asinh(\beta_0/c)$ and $\phi''(\beta)=1/\sqrt{\beta^2+c^2}$.

\step1 \emph{A partition function whose low-temperature limit is the minimiser.} For an inverse temperature $\kappa$,
\[
Z=\int d\beta\ e^{-\kappa\sum_j\phi_j(\beta_j)}\prod_i\delta\big(x_i^\top w-\sigma\xi_i\big).
\]
As $\kappa\to\infty$ the weight concentrates on $\hat\beta$. Averages over this weight follow from $\E\log Z$ (expectation over $X,\xi,\beta_0,b$).
\step2 \emph{Replica trick.} $\E\log Z=\lim_{n\to0}(\E Z^n-1)/n$. For integer $n$, $Z^n$ is an integral over $n$ copies $\beta^1,\dots,\beta^n$ (replicas), $w^a=\beta^a-b$.
\step3 \emph{Average over the inputs.} For fixed replicas the numbers $x_i^\top w^a$, $a=1..n$, are jointly Gaussian with mean zero and covariance
\[
Q_{ab}=\frac1d\sum_jw^a_jw^b_j .
\]
So $\E_{x_i}\prod_a\delta(x_i^\top w^a-\sigma\xi_i)=\mathcal N(\sigma\xi_i\mathbf 1;0,Q)$, the density of $\mathcal N(0,Q)$ at the point $\sigma\xi_i\mathbf 1$. The samples are independent:
\[
\E Z^n=\E\int\prod_ad\beta^a\ e^{-\kappa\sum_{a,j}\phi_j(\beta^a_j)}\ \Big[\E_\xi\,\mathcal N(\sigma\xi\mathbf 1;0,Q)\Big]^{\alpha d}.
\]
\step4 \emph{Make $Q$ an integration variable.} Insert $1=\int\prod_{a\le b}dQ_{ab}\,\delta(dQ_{ab}-\sum_jw^a_jw^b_j)$ and write each $\delta$ as a Fourier integral with a conjugate variable $\hat Q_{ab}$. Up to constants,
\[
\E Z^n=\int dQ\,d\hat Q\ \exp\,d\Big\{\tfrac12\sum_{ab}\hat Q_{ab}Q_{ab}+\log\E\!\int\prod_ad\beta^a\,e^{-\kappa\sum_a\phi(\beta^a)-\frac12\sum_{ab}\hat Q_{ab}w^aw^b}+\alpha\log\E_\xi\mathcal N(\sigma\xi\mathbf 1;0,Q)\Big\}.
\]
The coordinates have decoupled: the middle term is an integral over the $n$ replicas of \emph{one} coordinate, $\E$ over its $(\beta_0,b)$. For large $d$ the integral over $Q,\hat Q$ is given by the stationary point of the bracket.
\step5 \emph{Replica-symmetric ansatz.} All replicas are equivalent:
\[
Q_{aa}=q+\Delta,\quad Q_{ab}=q\ (a\neq b);\qquad\hat Q_{aa}=\hat A-\hat r,\quad\hat Q_{ab}=-\hat r\ (a\neq b).
\]
$q$ is the overlap of two different replicas, $\Delta$ the extra self-overlap (thermal spread). We now evaluate the three terms to first order in $n$.
\step6 \emph{First term.} $\frac12\sum_{ab}\hat Q_{ab}Q_{ab}=\frac n2(\hat A-\hat r)(q+\Delta)-\frac{n(n-1)}2\hat rq\approx\frac n2\big[\hat A(q+\Delta)-\hat r\Delta\big]$.
\step7 \emph{Second term.} $-\frac12\sum_{ab}\hat Q_{ab}w^aw^b=-\frac{\hat A}2\sum_a(w^a)^2+\frac{\hat r}2\big(\sum_aw^a\big)^2$. Linearise the square with a Gaussian $z\sim\mathcal N(0,1)$: $e^{\frac{\hat r}2s^2}=\E_ze^{\sqrt{\hat r}zs}$. The replicas become independent given $z$:
\[
\log\E\Big[\int d\beta\,e^{-\kappa\phi(\beta)-\frac{\hat A}2w^2+\sqrt{\hat r}zw}\Big]^n\approx n\,\E\log\int d\beta\,e^{-\kappa\phi(\beta)-\frac{\hat A}2w^2+\sqrt{\hat r}zw} .
\]
\step8 \emph{Third term.} $Q=\Delta I+q\mathbf 1\mathbf 1^\top$ has $\det Q=\Delta^{n-1}(\Delta+nq)$ and $\mathbf 1^\top Q^{-1}\mathbf 1=n/(\Delta+nq)$. Hence
\[
\log\mathcal N(\sigma\xi\mathbf 1;0,Q)=-\tfrac n2\log2\pi-\tfrac{n-1}2\log\Delta-\tfrac12\log(\Delta+nq)-\frac{\sigma^2\xi^2n}{2(\Delta+nq)}\approx-\frac n2\Big[\log(2\pi\Delta)+\frac{q+\sigma^2\xi^2}{\Delta}\Big],
\]
and $\E\xi^2=1$.
\step9 \emph{The function to make stationary} ($\frac1{nd}\log\E Z^n$):
\[
\Phi=\tfrac12\big[\hat A(q+\Delta)-\hat r\Delta\big]+\E\log\!\int d\beta\,e^{-\kappa\phi(\beta)-\frac{\hat A}2w^2+\sqrt{\hat r}zw}-\frac\alpha2\Big[\log(2\pi\Delta)+\frac{q+\sigma^2}{\Delta}\Big].
\]
\step{10} \emph{Stationarity.} Let $\langle\cdot\rangle$ be the average over $\beta$ with the weight in the integral.
\begin{align*}
\partial_q:&\quad\hat A=\alpha/\Delta, &\partial_{\hat A}:&\quad q+\Delta=\E\langle w^2\rangle,\\
\partial_\Delta:&\quad\hat r=\alpha(q+\sigma^2)/\Delta^2\ \ (\text{using }\hat A=\alpha/\Delta), &\partial_{\hat r}:&\quad\Delta=\E\big[\langle w^2\rangle-\langle w\rangle^2\big].
\end{align*}
(The last uses $\partial_{\hat r}\E f(\sqrt{\hat r}z)=\frac12\E f''$ for Gaussian $z$.) So $q=\E\langle w\rangle^2$.
\step{11} \emph{Zero temperature.} As $\kappa\to\infty$ the spread $\Delta$ vanishes like $1/\kappa$. Set $\Delta=\chi/\kappa$, $\hat A=\kappa A$, $\hat r=\kappa^2r$, $\eta=\sqrt rz$. The weight becomes $e^{-\kappa[\phi(\beta)+\frac A2w^2-\eta w]}$ and concentrates at the minimiser of the bracket: $\phi'(\hat\beta)+A(\hat\beta-b)=\eta$, which is \eqref{ch1}. Its spread is $1/\kappa(\phi''+A)$, so $\chi=\E\,1/(\phi''(\hat\beta)+A)$; differentiating \eqref{ch1} in $\eta$ shows this equals $\E\,\partial\hat\beta/\partial\eta$. The first three relations of Step 10 give $A=\alpha/\chi$, $r=\alpha(q+\sigma^2)/\chi^2$, $q=\E(\hat\beta-b)^2$.

\textbf{Reading.} $\eta$ is the field a coordinate feels from the fitted samples; $\chi$ is how much the coordinate gives way to it; $A$ is the stiffness with which the samples hold it at the target.

\textbf{Pretraining} ($\beta_0=0$, $b=\bs$): $\asinh(\beta/c)+A(\beta-\bs)=\eta$. Dividing by $A$: $\beta=\arg\min_x\{\frac1A\psi_c(x)+\frac12(x-\bs-\eta/A)^2\}$ with noise variance $r/A^2=(q+\sigma^2)/\alpha$, and $\chi A=\alpha$ reads $\E\,\partial\beta/\partial(\bs+\eta/A)=\alpha$.

\textbf{Weight decay kept on} (Result 5). The constraints are replaced by the factor $e^{-\frac\kappa2\sum_i(x_i^\top w-\sigma\xi_i)^2}$ and $\phi$ by $\frac\lambda2|\beta|$. Step 8 becomes a Gaussian integral: $\E_p\,e^{-\frac\kappa2\sum_a(p_a-\sigma\xi)^2}$ with $p\sim\mathcal N(0,Q)$, which to first order in $n$ is $\exp-\frac n2\big[\log(1+\kappa\Delta)+\frac{\kappa(q+\sigma^2\xi^2)}{1+\kappa\Delta}\big]$. With $\Delta=\chi/\kappa$ the same steps give
\[
A=\frac{\alpha}{1+\chi},\qquad r=\frac{\alpha(q+\sigma^2)}{(1+\chi)^2},\qquad\hat\beta=\operatorname{soft}\Big(b+\frac\eta A,\ \frac\lambda{2A}\Big).
\]

\section{Replica method: blocks of coordinates and groups of samples}
Samples come in groups $g$ ($\alpha_gd$ samples) that see only some blocks $B$ of coordinates (block $B$ is a fraction $p_B$ of the coordinates), with targets generated by $b^g$ and noise $\sigma_g$.

\textbf{Result 14.} A coordinate of block $B$ obeys
\begin{equation}\label{chB}
\phi'(\hat\beta)+\sum_{g\ni B}A_g(\hat\beta-b^g)=\eta_B,\qquad\operatorname{Var}\eta_B=\sum_{g\ni B}r_g ,
\end{equation}
\[
A_g=\frac{\alpha_g}{\chi_g},\qquad r_g=\frac{\alpha_g(q_g+\sigma_g^2)}{\chi_g^2},\qquad q_g=\sum_{B\in g}p_B\,\E_B(\hat\beta-b^g)^2,\qquad\chi_g=\sum_{B\in g}p_B\,\E_B\frac{\partial\hat\beta}{\partial\eta_B},
\]
where $g\ni B$ means that group $g$ sees block $B$.

\emph{Derivation: the changes to Section~\ref{rep1}.} Step 3: for a sample of group $g$, $x_i^\top(\beta^a-b^g)$ has covariance $Q^g_{ab}=\frac1d\sum_{j\in\text{blocks of }g}(\beta^a_j-b^g_j)(\beta^b_j-b^g_j)$, one overlap matrix per group. Step 4: one conjugate $\hat Q^g$ per group; the energetic term is $\sum_g\alpha_g\log\mathcal N(\cdot;0,Q^g)$, and a coordinate of block $B$ carries $-\frac12\sum_{g\ni B}\sum_{ab}\hat Q^g_{ab}(\beta^a-b^g)(\beta^b-b^g)$, with weight $p_B$ in the entropic term. Steps 5--11 go through group by group: each group has its own $(q_g,\chi_g,A_g,r_g)$ and its own independent Gaussian $z_g$; a coordinate of block $B$ sums the contributions of the groups that see it.

\section{Replica method: finetuning on samples used in pretraining}\label{rep2}
Now the start of stage 2 depends on the data, so Section~\ref{rep1} does not apply. Take first one block and one group: pretraining from zero on all $\alpha d$ samples with labels $x_i^\top\bs$; then rescaling by $S$; then finetuning on $\alpha_Fd$ \emph{of the same samples} to targets $b\,y_i$. Stage-2 cost per coordinate: $\phi_2(\beta_2;\beta_1)=D_{Sc}(\beta_2,S\beta_1)$, so $\phi_2'=\asinh(\beta_2/Sc)-\asinh(\beta_1/c)$.

\textbf{Result 15.} With $u=\beta_1-\bs$ and $w=\beta_2-b\bs$, one coordinate obeys
\begin{align}
\text{stage 1:}&\quad\asinh(\beta_1/c)+A_1u=\eta_1,\label{s1}\\
\text{stage 2:}&\quad\asinh(\beta_2/Sc)-\asinh(\beta_1/c)+A_2\,(w-\varrho\,u)=\eta_2,\label{s2}
\end{align}
with $(\eta_1,\eta_2)$ jointly Gaussian,
\[
A_1=\frac\alpha{\chi_1},\quad\operatorname{Var}\eta_1=\frac{\alpha\,\E u^2}{\chi_1^2},\qquad
A_2=\frac{\alpha_F}{\chi_2},\quad\operatorname{Var}\eta_2=\frac{\alpha_F\,\E(w-\varrho u)^2}{\chi_2^2},\quad\operatorname{Cov}(\eta_1,\eta_2)=\frac{\alpha_F\,\E[u\,(w-\varrho u)]}{\chi_1\chi_2},
\]
\[
\chi_1=\E\frac{\partial\beta_1}{\partial\eta_1},\qquad\chi_2=\E\frac{\partial\beta_2}{\partial\eta_2},\qquad\varrho=\frac1{\chi_1}\E\frac{d\beta_2}{d\eta_1}\quad\text{(through $\beta_1$, at fixed $\eta_2$)}.
\]
\emph{Reading.} A reused sample was already fitted in stage 1. Its stage-2 residual is not $w$ but $w-\varrho u$: the part of the new misfit that is not explained by the misfit it already absorbed. $\varrho$ measures how much of a change of the pretrained model is still present after finetuning. For fresh samples $\varrho$ is replaced by $0$ and $\operatorname{Cov}(\eta_1,\eta_2)=0$.

\step1 \emph{What to average.} The pretrained model is a sample from the stage-1 weight
\[e^{-\kappa_1\sum_j\psi_c(\beta_{1,j})}\prod_i\delta(x_i^\top u)\big/Z_1 ,\]
and we need $\E\log Z_2(\beta_1)$, with $Z_2$ the stage-2 partition function started at $\beta_1$. Write $1/Z_1=\lim_{m\to0}Z_1^{m-1}$ and $\log Z_2=\lim_{n\to0}\partial_nZ_2^n$. This gives $m$ replicas of stage 1, $u^1,\dots,u^m$, and $n$ replicas of stage 2, $w^1,\dots,w^n$, all started from the \emph{first} stage-1 replica $u^1$.
\step2 \emph{Average over the inputs.} A sample used in both stages imposes $x^\top u^a=0$ for all $a$ and $x^\top w^c=0$ for all $c$. These $m+n$ numbers are jointly Gaussian with covariance
\[
\mathcal Q=\begin{pmatrix}Q_1&C\\ C^\top&Q_2\end{pmatrix},\qquad Q_1=\tfrac1d\textstyle\sum_ju^au^b,\quad Q_2=\tfrac1d\sum_jw^cw^{c'},\quad C_{ac}=\tfrac1d\sum_ju^aw^c ,
\]
and the sample contributes the density $\mathcal N(0;0,\mathcal Q)=(2\pi)^{-(m+n)/2}(\det\mathcal Q)^{-1/2}$. A sample used only in stage 1 contributes $\mathcal N(0;0,Q_1)$.
\step3 \emph{Ansatz.} $Q_1$ and $Q_2$ replica symmetric with $(q_1,\Delta_1)$ and $(q_2,\Delta_2)$. Cross overlaps: $C_{ac}=q_{12}$ for $a\neq1$, and $C_{1c}=q_{12}+\Delta_{12}$: a stage-2 replica overlaps more with the replica it started from.
\step4 \emph{The determinant.} $\det\mathcal Q=\det Q_1\cdot\det(Q_2-C^\top Q_1^{-1}C)$. Here $C=\mathbf c\,\mathbf 1_n^\top$ with $\mathbf c=q_{12}\mathbf 1_m+\Delta_{12}\mathbf e_1$, and $Q_1^{-1}=\frac1{\Delta_1}\big[I-\frac{q_1}{\Delta_1+mq_1}\mathbf 1\mathbf 1^\top\big]$. So $C^\top Q_1^{-1}C=K\,\mathbf 1_n\mathbf 1_n^\top$ with, as $m\to0$,
\[
K=\mathbf c^\top Q_1^{-1}\mathbf c=\frac1{\Delta_1}\Big[2q_{12}\Delta_{12}+\Delta_{12}^2-\frac{q_1}{\Delta_1}\Delta_{12}^2\Big].
\]
Then $Q_2-K\mathbf 1\mathbf 1^\top$ is replica symmetric with $q_2\to q_2-K$, and by Step 8 of Section~\ref{rep1} the part of $\log\mathcal N(0;0,\mathcal Q)$ proportional to $n$ is $-\frac n2\big[\log(2\pi\Delta_2)+\frac{q_2-K}{\Delta_2}\big]$. ($q_2-K$ is the variance of $x^\top w$ given that $x^\top u^a=0$.) The factor $\det Q_1$ does not involve stage 2: as $n\to0$ stage 1 is unaffected and obeys Result 13.
\step5 \emph{Conjugate variables.} For $Q_2$: $\hat A_2,\hat r_2$ as before. For the cross overlaps: $-\hat r_{12}$ for every pair $(a,c)$ plus $\hat A_{12}$ for the pairs $(1,c)$. The coupling term, to order $n$ and as $m\to0$:
\[
\tfrac12\big[\hat A_2(q_2+\Delta_2)-\hat r_2\Delta_2\big]+\hat A_{12}(q_{12}+\Delta_{12})-\hat r_{12}\Delta_{12}.
\]
(The cross sum is $n[-\hat r_{12}(mq_{12}+\Delta_{12})+\hat A_{12}(q_{12}+\Delta_{12})]$.)
\step6 \emph{One coordinate.} The exponent is
\[
-\tfrac{\hat A_1}2\sum_a(u^a)^2-\tfrac{\hat A_2}2\sum_c(w^c)^2-\hat A_{12}\,u^1\sum_cw^c+\tfrac{\hat r_1}2\Big(\sum_au^a\Big)^2+\hat r_{12}\Big(\sum_au^a\Big)\Big(\sum_cw^c\Big)+\tfrac{\hat r_2}2\Big(\sum_cw^c\Big)^2 .
\]
The last three terms are a quadratic form in $(\sum u,\sum w)$; linearise it with a Gaussian pair $(z_1,z_2)$ of covariance $\begin{pmatrix}\hat r_1&\hat r_{12}\\\hat r_{12}&\hat r_2\end{pmatrix}$: it becomes $\E\,e^{z_1\sum_au^a+z_2\sum_cw^c}$. Given $(z_1,z_2)$ the replicas factorise:
\[
\Big[\int du\,e^{-\kappa_1\psi_c-\frac{\hat A_1}2u^2+z_1u}\Big]^{m-1}\int du^1\,e^{-\kappa_1\psi_c-\frac{\hat A_1}2(u^1)^2+z_1u^1}\Big[\int dw\,e^{-\kappa_2\phi_2(\beta_2;\beta_1)-\frac{\hat A_2}2w^2-\hat A_{12}u^1w+z_2w}\Big]^n .
\]
As $m\to0$ the first bracket becomes $1/Z_1(z_1)$, which normalises the integral over $u^1$: $u^1$ is a sample from the stage-1 single-coordinate weight at field $z_1$. To order $n$ the term is $n\,\E\big\langle\log Z_2(z_2,u_1)\big\rangle_{u_1|z_1}$ with $Z_2$ the last bracket.
\step7 \emph{The function to make stationary} (order $n$):
\[
\Phi_2=\tfrac12\big[\hat A_2(q_2+\Delta_2)-\hat r_2\Delta_2\big]+\hat A_{12}(q_{12}+\Delta_{12})-\hat r_{12}\Delta_{12}+\E\big\langle\log Z_2\big\rangle-\frac{\alpha_F}2\Big[\log(2\pi\Delta_2)+\frac{q_2-K}{\Delta_2}\Big].
\]
\step8 \emph{Stationarity in the overlaps.}
\begin{align*}
\partial_{q_2}:&\ \hat A_2=\frac{\alpha_F}{\Delta_2}, &\partial_{\Delta_2}:&\ \hat r_2=\frac{\alpha_F(q_2-K)}{\Delta_2^2},\\
\partial_{q_{12}}:&\ \hat A_{12}=-\frac{\alpha_F}{2\Delta_2}\frac{\partial K}{\partial q_{12}}=-\frac{\alpha_F\Delta_{12}}{\Delta_1\Delta_2}, &\partial_{\Delta_{12}}:&\ \hat r_{12}=\hat A_{12}+\frac{\alpha_F}{2\Delta_2}\frac{\partial K}{\partial\Delta_{12}}=\frac{\alpha_F}{\Delta_1\Delta_2}\Big[q_{12}-\frac{q_1}{\Delta_1}\Delta_{12}\Big].
\end{align*}
Stationarity in the conjugates identifies the overlaps with averages: $q_2+\Delta_2=\E\langle w^2\rangle$, $\Delta_2=\E[\langle w^2\rangle-\langle w\rangle^2]$, $q_{12}+\Delta_{12}=\E\langle u_1w\rangle$, $\Delta_{12}=\E[\langle u_1w\rangle-\langle u_1\rangle\langle w\rangle]$.
\step9 \emph{Zero temperature.} $\Delta_1=\chi_1/\kappa_1$, $\Delta_2=\chi_2/\kappa_2$. $\Delta_{12}$ is the joint thermal fluctuation of $u_1$ and $w$; $u_1$ fluctuates by $O(1/\kappa_1)$ and $w$ follows it, so $\Delta_{12}=\chi_{12}/\kappa_1$ with $\chi_{12}=\E\,d\beta_2/d\eta_1$. Put $\varrho=\chi_{12}/\chi_1$. Then
\[
K\to2\varrho q_{12}-\varrho^2q_1,\qquad q_2-K=\E(w-\varrho u)^2,
\]
\[
\hat A_2=\kappa_2\frac{\alpha_F}{\chi_2},\qquad\hat A_{12}=-\kappa_2\frac{\alpha_F\varrho}{\chi_2},\qquad\hat r_2=\kappa_2^2\frac{\alpha_F\E(w-\varrho u)^2}{\chi_2^2},\qquad\hat r_{12}=\kappa_1\kappa_2\frac{\alpha_F\E[u(w-\varrho u)]}{\chi_1\chi_2}.
\]
With $z_s=\kappa_s\eta_s$ the stage-2 weight is $e^{-\kappa_2[\phi_2+\frac{A_2}2w^2-A_2\varrho\,u_1w-\eta_2w]}$; its minimiser satisfies \eqref{s2}, and the covariances of $(\eta_1,\eta_2)$ are $\hat r/\kappa\kappa$.

\textbf{Check.} Solving \eqref{s1}--\eqref{s2} directly ($\alpha=0.4$, $\alpha_F=0.2$, $b=0$, $S=1$, $k=10^{-2}$) gives $\E\beta_2^2=0.0721$, $\E(\beta_2-\bs)^2=0.855$, $\E(\beta_2-\beta_1)^2=0.322$; the cavity solver used for the table below gives $0.0721$, $0.852$, $0.323$. At $k=10^{-5}$, $\alpha=0.3$, $\alpha_F=0.15$: $0.0248$, $0.936$, $0.372$ against $0.0249$, $0.933$, $0.369$.

\textbf{Several groups and blocks.} As in Result 14, each group $g$ used in stage 2 has its own target multiple $b_g$, residual $e_g=(\beta_2-b_g\bs)-\varrho_g(\beta_1-\bs)$ and $\varrho_g=\chi_{12}^g/\chi_1^g$ (zero for fresh samples); the $\chi$ of a group average over the blocks it sees. For a coordinate of block $B$:
\begin{equation}\label{s2g}
\asinh(\beta_2/Sc)-\asinh(\beta_1/c)+\sum_{g\ni B}\frac{\alpha_g}{\chi_2^g}\,e_g=\eta_2,\quad
\operatorname{Var}\eta_2=\sum_{g\ni B}\frac{\alpha_g\E_ge_g^2}{(\chi_2^g)^2},\quad
\operatorname{Cov}(\eta_1,\eta_2)=\sum_{g\ni B}\frac{\alpha_g\E_g[u\,e_g]}{\chi_1^g\chi_2^g},
\end{equation}
where $\E_g$ averages over the blocks group $g$ sees, weighted by $p_B$. The retain term is a group with $b=1$; forget-to-zero is $b=0$; a negated target is $b<0$.

\section{Replica method: the references and the metrics}
\textbf{References.} The model retrained without $F$ and the retain-only model are each a single stage from zero on fewer groups: Result 14. Their distance to the finetuned model needs their overlap with it.

\textbf{Result 16.} Two models $s,s'$ trained on overlapping data have correlated fields: for a coordinate of block $B$,
\[
\operatorname{Cov}(\eta_s,\eta_{s'})=\sum_{g\ni B\text{ used by both}}\frac{\alpha_g\,\E_g[e^s_ge^{s'}_g]}{\chi^g_s\chi^g_{s'}},
\]
with $e^s_g$ the residual of model $s$ on group $g$ ($\beta_s-b_g\bs$ for a stage from zero, $e_g$ of \eqref{s2g} for stage 2). All overlaps follow by solving the coordinate equations with these jointly Gaussian fields.

\emph{Derivation.} Replicate both models ($m$ and $n$ replicas). A shared sample constrains both, so Step 4 of Section~\ref{rep2} applies with a cross block $q_{ss'}\mathbf 1\mathbf 1^\top$ (no distinguished replica). Then $K=q_{ss'}^2\,\mathbf 1^\top Q_s^{-1}\mathbf 1\approx mq_{ss'}^2/\Delta_s$, and the density gains $\frac{mn}2\,q_{ss'}^2/\Delta_s\Delta_{s'}$ per shared sample. Stationarity in $q_{ss'}$ sets the conjugate, which is the covariance of the two linearising Gaussians (Step 6), to $\alpha_gq_{ss'}/\Delta_s\Delta_{s'}$; at zero temperature this is the formula. For stage 2, $q_2\to q_2-K$ replaces $w$ by $w-\varrho u$ throughout. Result 15's $\operatorname{Cov}(\eta_1,\eta_2)$ is the case $s=1$, $s'=2$.

\textbf{Metrics.} With $\E_B$ the average over block $B$:
\begin{itemize}[nosep,leftmargin=1.2em]
\item error on a class's fresh inputs: $\sum_{B\text{ seen}}p_B\E_B(\beta_2-\bs)^2$;
\item error on $F$ after convergence: $(1-b)^2$ times the label variance (Result 11);
\item error on a left-over sample, which was fitted in stage 1 and is unused in stage 2: $\sum_{B\text{ seen}}p_B\,\E_B\big[(\beta_2-\bs)-\varrho(\beta_1-\bs)\big]^2$, by the same conditioning as in Step 4;
\item distance to a reference: $\sum_Bp_B\E_B(\beta_2-\beta_{\rm ref})^2$.
\end{itemize}

\textbf{Check against simulation} (end of finetuning, $t=40$, $d=3000$; the simulation is still drifting slowly upward with shared coordinates). Distance to the retrained model, distance to the retain-only model:
\begin{center}\small\begin{tabular}{@{}llcccc@{}}\toprule
blocks & & forget-to-zero & negated $-0.5$ & $S=0.3$ & $S=3$\\\midrule
$(1,0,1)$ & equations & $0.109,\ 0.012$ & & $0.132,\ 0.002$ & $0.113,\ 0.109$\\
& simulation & $0.113,\ 0.012$ & & $0.135,\ 0.002$ & $0.107,\ 0.096$\\
$(1,0.4,1)$ & equations & $0.149,\ 0.110$ & $0.281,\ 0.271$ & $0.168,\ 0.083$ & $0.170,\ 0.220$\\
& simulation & $0.140,\ 0.096$ & $0.260,\ 0.249$ & $0.159,\ 0.071$ & $0.152,\ 0.193$\\\bottomrule
\end{tabular}\end{center}

\section{The same equations by the cavity method}
A second derivation, with different assumptions (removing one row or column of $X$ changes the rest to first order only; sums of many independent terms equal their means). It also explains the quantities: with $V=-\chi$ and $a_i$ the accumulated gradient of sample $i$ (Result 2, Step 2), the end point of pretraining satisfies
\begin{equation}\label{kkt}
\theta_j=-\sum_iX_{ji}a_i,\qquad\beta_j=c\sinh\theta_j,\qquad\sum_jX_{ji}\beta_j=y_i .
\end{equation}
\textbf{Result 17 (Result 13 again).} One coordinate and one sample obey
\begin{equation}\label{st}
\theta+A(\beta-\bs)=-\eta,\quad\beta=c\sinh\theta;\qquad\zeta+Va=y;\qquad A=-\frac\alpha V,\quad V=\E\frac{\partial\beta}{\partial\eta}=-\E\frac{\sqrt{\beta^2+c^2}}{1+A\sqrt{\beta^2+c^2}} ,
\end{equation}
where $\eta\sim\mathcal N(0,\alpha\,\E a^2)$ and $(z,\zeta)$, with $y=z$, is Gaussian with the covariance of $(\bs,\beta)$ over coordinates.

\step1 \emph{Add a sample} (index $0$) to a system solved without it. Its inputs are independent of that solution $\beta^{\setminus0}$, so $\zeta=x_0^\top\beta^{\setminus0}$ and $z=x_0^\top\bs$ are Gaussian with variances $\E\beta^2$, $\E(\bs)^2$ and covariance $\E\beta\bs$. Once added, the sample gets some $a_0$, which by \eqref{kkt} shifts every $\theta_j$ by $-X_{j0}a_0$. Call $\eta_j$ the quantity with $\theta_j=-\eta_j+\dots$; it shifts by $X_{j0}a_0$, so $\beta_j=\beta_j^{\setminus0}+\frac{\partial\beta_j}{\partial\eta_j}X_{j0}a_0$. The sample's prediction is
\[
\sum_jX_{j0}\beta_j=\zeta+\Big(\sum_jX_{j0}^2\frac{\partial\beta_j}{\partial\eta_j}\Big)a_0=\zeta+Va_0,
\]
using $\sum_jX_{j0}^2(\cdot)\to\E(\cdot)$. Interpolation gives $\zeta+Va_0=y$, so $a_0=(y-\zeta)/V$.
\step2 \emph{Add a coordinate} (index $0$). Without it the samples have $a_i^{\setminus0}$, independent of the row $X_{0\cdot}$. With it, sample $i$'s prediction gains $X_{0i}\beta_0$ and its label gains $X_{0i}\bs_0$: the rest of the model must now fit $y_i-X_{0i}(\beta_0-\bs_0)$. By Step 1, $a_i$ responds to its label with slope $1/V$: $a_i=a_i^{\setminus0}-\frac1VX_{0i}(\beta_0-\bs_0)$. Insert in \eqref{kkt}:
\[
\theta_0=-\sum_iX_{0i}a_i^{\setminus0}+\frac1V\Big(\sum_iX_{0i}^2\Big)(\beta_0-\bs_0)=-\eta_0-A(\beta_0-\bs_0),
\]
with $\sum_iX_{0i}^2\to\alpha$, $A=-\alpha/V$, and $\eta_0=\sum_iX_{0i}a_i^{\setminus0}$ Gaussian with variance $\frac1d\sum_ia_i^2=\alpha\,\E a^2$.
\step3 \emph{Close.} Differentiate the coordinate equation in $\eta$: $d\theta\,(1+A\sqrt{\beta^2+c^2})=-d\eta$, and $d\beta=\sqrt{\beta^2+c^2}\,d\theta$. This gives $V$.

This is Result 13 with $\eta\to-\eta$, $V=-\chi$, $q+\sigma^2=\operatorname{Var}(y-\zeta)$. The two-stage equations follow in the same way: a reused sample has $\zeta_2+V_{22}a_2+V_{21}a_1=by$ with $V_{21}=\E\,d\beta_2/d\eta_1$, so $a_2=\big[(by-\zeta_2)-\varrho(y-\zeta_1)\big]/V_{22}$, which is the residual $w-\varrho u$ of Result 15.

\section{Large $d$: the finetuning trajectory (cavity method in time)}
The replica method gives end points only. For trajectories the same cavity steps are applied at every time step.

Finetuning is run in steps of size $\delta$: $u\leftarrow u(1-2\delta G_n-\delta\lambda)$, $v\leftarrow v(1+2\delta G_n-\delta\lambda)$, where $G_{n,j}=\sum_iX_{ji}g_{n,i}$ and $g_{n,i}=w_i(p_{n,i}-t_i)$ is the gradient of sample $i$ at step $n$, $p_{n,i}$ its prediction. The pretrained model is called step $0$: its ``gradient'' is $g_0=a$ and its field is $\eta_0$ of \eqref{st}.

\textbf{Result 18.} One sample and one coordinate obey, for $n\ge1$,
\begin{equation}\label{dyn}
p_n=\zeta_n+\sum_{l<n}R_\beta(n,l)\,g_l,\qquad G_n=\eta_n+\sum_{l\le n}R_g(n,l)\,\beta_l+\mu_n\bs ,
\end{equation}
with $R_\beta(n,l)=\E\frac{\partial\beta_n}{\partial G_l}$, $R_g(n,l)=\sum_a\alpha_a\E_a\frac{\partial g_n}{\partial\zeta_l}$, $\mu_n=\sum_a\alpha_a\E_a\frac{\partial g_n}{\partial z}$; $(z,\zeta_0,\zeta_1,\dots)$ Gaussian with the covariances of $(\bs,\beta_0,\beta_1,\dots)$; $(\eta_0,\eta_1,\dots)$ Gaussian with covariances $\sum_a\alpha_a\E_a\,g_ng_m$.

\step1 \emph{Add a sample.} Without it the weights at every step are independent of its inputs, so $\zeta_n=x_0^\top\beta_n^{\setminus0}$ are jointly Gaussian with the stated covariances. With it, the field of coordinate $j$ at each earlier step $l$ gains $X_{j0}g_{l,0}$, so $\beta_{n,j}=\beta_{n,j}^{\setminus0}+\sum_{l<n}\frac{\partial\beta_{n,j}}{\partial G_{l,j}}X_{j0}g_{l,0}$. Multiply by $X_{j0}$ and sum over $j$: the first equation of \eqref{dyn}.
\step2 \emph{Add a coordinate.} Sample $i$'s predictions-without-it at all steps $l\le n$ gain $X_{0i}\beta_{l,0}$, and its label input gains $X_{0i}\bs_0$. To first order $g_{n,i}=g_{n,i}^{\setminus0}+\sum_{l\le n}\frac{\partial g_{n,i}}{\partial\zeta_{l,i}}X_{0i}\beta_{l,0}+\frac{\partial g_{n,i}}{\partial z_i}X_{0i}\bs_0$. Multiply by $X_{0i}$ and sum over $i$: the second equation, with $\eta_n=\sum_iX_{0i}g_{n,i}^{\setminus0}$.
\step3 \emph{The responses} follow by differentiating the recursions. Samples: from \eqref{dyn} and $g_n=w(p_n-t)$, with $t$ containing $b$ times the label,
\[
\frac{\partial g_n}{\partial\zeta_m}=w\Big(\mathbf 1_{n=m}+\sum_{l<n}R_\beta(n,l)\frac{\partial g_l}{\partial\zeta_m}\Big),\qquad\frac{\partial g_n}{\partial z}=w\Big(\sum_{l<n}R_\beta(n,l)\frac{\partial g_l}{\partial z}-b\Big).
\]
Coordinates: differentiate the update, $du\leftarrow du\,(1-2\delta G_n-\delta\lambda)-2\delta u\,dG_n$ (and $v$ with the opposite sign of $G$), where $dG_n=\mathbf 1_{n=m}+\sum_{l\le n}R_g(n,l)\,d\beta_l$ for a perturbation of the field at step $m$, and $d\beta=2(u\,du-v\,dv)$.
\step4 \emph{Step 0.} From \eqref{st}: $\partial g_0/\partial\zeta_0=-1/V$, $\partial g_0/\partial z=+1/V$, $\partial\beta_0/\partial\eta_0$ as in Result 17, and the starting $u,v$ are $\sqrt{c/2}\,e^{\pm\theta/2}$ times $\sqrt S$.

The four metrics: errors on the two distributions are $\sum_Bq_B\E_B(\beta_n-\bs)^2$ over the blocks each class sees; errors on $F$ and on $F_{\rm full}\setminus F$ are $\E(p_n-y)^2$ over those groups (the left-over samples have $g_n=0$ for $n\ge1$, so $p_n=\zeta_n+R_\beta(n,0)g_0$).

\section{How the equations are solved}
\textbf{Idea.} Represent the Gaussians by two independent random matrices on $N$ sampled coordinates and $\alpha N$ sampled samples: $\zeta_n=G_1\beta_n$ and $\eta_n=G_2g_n$. They have the required covariances automatically. Because $G_1$ and $G_2$ are independent, no self-interaction arises; it is added explicitly by the $R$ terms of \eqref{dyn}. The expectations defining $R_\beta,R_g,\mu$ are averages over the sampled paths.

\textbf{Starting from the pretrained model.} \eqref{st} is solved first, with fixed standard normal draws, giving vectors $\beta_0,a,\zeta_0,\eta_0,z$. The matrices must reproduce them: $G_1\beta_0=\zeta_0$, $G_1\bs=z$, $G_2a=\eta_0$.

\textbf{Result 19.} If $\tilde G$ has independent Gaussian entries, then $G=\tilde G-(\tilde GB-Z)(B^\top B)^{-1}B^\top$ has the law of $\tilde G$ conditioned on $\tilde GB=Z$. \emph{Derivation:} $\tilde G(I-P)$ with $P=B(B^\top B)^{-1}B^\top$ is independent of $\tilde GB$ (Gaussian, zero covariance), and $G=\tilde G(I-P)+Z(B^\top B)^{-1}B^\top$ keeps the first and replaces $\tilde GP$ by its required value.

Then \eqref{dyn} is run forward once. Cost grows as (steps)$^3$, which limits theory curves to $t\le6$.

\section{What is and is not derived}
\begin{center}\small\begin{tabular}{@{}p{7.5cm}p{7.5cm}@{}}\toprule
derived & not derived\\\midrule
Results 1--12: exact & why the distance to the retrained model never falls below about half\\
Results 13--16: replica, assuming replica symmetry; end points agree with simulation (table in Section 11) & a closed formula for the spread in the rich regime beyond the first-order rule \eqref{lin} (about $1.3\times$ too high at large $|F|$)\\
Results 17--19: cavity; trajectories agree with simulation ($94\%$ of 2868 comparisons within 2 s.e.) & finetuning after SGD pretraining; end points when the targets cannot all be fitted\\\bottomrule
\end{tabular}\end{center}
\end{document}
